% Szendr\H{o}i, B., "Non-commutative Donaldson--Thomas theory and the conifold",
%   Geom. Topol. 12 (2008) 1171--1202 (arXiv:0705.3419).
% Van den Bergh, M., "Three-dimensional flops and non-commutative rings",
%   Duke Math. J. 122 (2004) 423--455 (arXiv:math/0207170).
% Bridgeland, T., "Flops and derived categories", Invent. Math. 147 (2002)
%   613--632 (arXiv:math/0009053).
% Morrison, A., Mozgovoy, S., Nagao, K., and Szendr\H{o}i, B.,
%   "Motivic Donaldson--Thomas invariants of the conifold and the refined
%   topological vertex" (arXiv:1107.5017).

\providecommand{\CoHA}{\mathrm{CoHA}}
\providecommand{\Hom}{\mathrm{Hom}}
\providecommand{\End}{\mathrm{End}}
\providecommand{\Spec}{\mathrm{Spec}}
\providecommand{\Coh}{\mathrm{Coh}}
\providecommand{\Mod}{\text{-}\mathrm{Mod}}
\providecommand{\modd}{\text{-}\mathrm{mod}}
\providecommand{\Db}{D^{b}}
\providecommand{\DT}{\mathrm{DT}}
\providecommand{\cO}{\mathcal O}
\providecommand{\bP}{\mathbb P}
\providecommand{\bC}{\mathbb C}

\section*{Szendr\H{o}i--Morita Correspondence for the Conifold}

\subsection*{The Node and Its Non-Commutative Resolution}

Let
\[
  R=\bC[x,y,z,w]/(xy-zw),\qquad X_0=\Spec R .
\]
This is the affine ordinary double point, a three-dimensional
Gorenstein normal singularity.  The small resolutions
\[
  Y_{\pm}\longrightarrow X_0
\]
contract one rational curve with normal bundle
\(\cO_{\bP^1}(-1)\oplus\cO_{\bP^1}(-1)\).  Van den Bergh's theorem on
threefold flops applies because the resolutions are crepant and have
one-dimensional fibres.  It supplies a tilting generator on each small
resolution.
For the conifold this generator may be written as
\(\cO_{Y_{\pm}}\oplus\cO_{Y_{\pm}}(\pm 1)\), and its endomorphism
algebra over the base is a non-commutative crepant resolution of \(R\).

The divisor class group of \(R\) is generated by the rank-one reflexive
ideal
\[
  I=(x,z)\subset R .
\]
The corresponding algebra is
\[
  \Lambda=\End_R(R\oplus I).
\]
In fractional-ideal form,
\[
  I^\vee=\Hom_R(I,R)=I^{-1}=(1,y/z)=(1,w/x),
\]
inside the fraction field of \(R\), and hence
\[
  \Lambda\cong
  \begin{pmatrix}
    R & I^\vee\\
    I & R
  \end{pmatrix}.
\]
The opposite small resolution replaces \(I\) by \(I^\vee\).  The two
NCCRs are derived Morita equivalent through the flop; after choosing the
standard involution of the node they are exchanged by an isomorphism of
the base.

\subsection*{The Klebanov--Witten Jacobi Algebra}

Let \(Q_{\mathrm{con}}\) be the quiver with vertices \(0,1\), arrows
\[
  a_1,a_2\colon 0\to 1,\qquad b_1,b_2\colon 1\to 0,
\]
and superpotential
\[
  W_{\mathrm{KW}}=
  a_1b_1a_2b_2-a_1b_2a_2b_1 .
\]
The Jacobi algebra is
\[
  J_{\mathrm{con}}=
  \bC Q_{\mathrm{con}}/
  \langle\partial_{a_i}W_{\mathrm{KW}},
         \partial_{b_j}W_{\mathrm{KW}}\rangle .
\]
The cyclic derivatives are
\[
\begin{aligned}
  \partial_{a_1}W_{\mathrm{KW}}&=b_1a_2b_2-b_2a_2b_1,\\
  \partial_{a_2}W_{\mathrm{KW}}&=b_2a_1b_1-b_1a_1b_2,\\
  \partial_{b_1}W_{\mathrm{KW}}&=a_2b_2a_1-a_1b_2a_2,\\
  \partial_{b_2}W_{\mathrm{KW}}&=a_1b_1a_2-a_2b_1a_1 .
\end{aligned}
\]

Szendr\H{o}i identifies \(J_{\mathrm{con}}\) with the conifold NCCR:
\[
  J_{\mathrm{con}}\cong\Lambda=\End_R(R\oplus I),
  \qquad
  Z(J_{\mathrm{con}})\cong R .
\]
One concrete normalization sends the arrows \(a_1,a_2\colon R\to I\) to
multiplication by \(x,z\), and sends \(b_1,b_2\colon I\to R\) to the two
maps \(1,y/z\in I^\vee\).  The length-two cycles then give
\[
  b_1a_1=x,\qquad b_1a_2=z,\qquad b_2a_1=w,\qquad b_2a_2=y,
\]
so the central relation is precisely
\[
  (b_2a_2)(b_1a_1)=(b_1a_2)(b_2a_1),
  \qquad xy=zw .
\]
The four Jacobi relations say that the two possible length-three paths
with fixed endpoints agree.  In the matrix algebra these identities are
exactly associativity of multiplication in the fractional ideals
\(I\) and \(I^\vee\), together with \(xy=zw\).  Thus the
Jacobi presentation and the endomorphism-ring presentation describe the
same \(R\)-algebra.

\subsection*{Morita and Derived Morita Statements}

There are two algebraic levels.

First, \(J_{\mathrm{con}}\cong\Lambda\) is an isomorphism of
\(R\)-algebras after the choice of generators above.  Therefore
\[
  J_{\mathrm{con}}\modd=\Lambda\modd
\]
as \(R\)-linear abelian categories in that normalization.

Second, Van den Bergh's tilting generator gives derived Morita
equivalences
\[
  \Db\Coh(Y_{\pm})\simeq\Db(\Lambda\modd).
\]
Composing either equivalence with \(J_{\mathrm{con}}\cong\Lambda\)
gives
\[
  \Db\Coh(Y_{\pm})\simeq\Db(J_{\mathrm{con}}\modd).
\]
Bridgeland's flop theorem identifies the two small resolutions through a
Fourier--Mukai equivalence.  The NCCR is the common non-commutative
pivot for this equivalence.

\subsection*{Donaldson--Thomas Chambers}

Szendr\H{o}i's non-commutative Donaldson--Thomas invariants count
framed cyclic \(J_{\mathrm{con}}\)-modules in the NCDT chamber.  The
commutative DT invariants count ideal sheaves on a chosen small
resolution.  The commutative DT chambers, PT chambers, the flop chamber,
and the NCDT chamber are different chambers in the stability space of
the framed conifold quiver.

Let \(M(q)=\prod_{n\geq 1}(1-q^n)^{-n}\).  In Szendr\H{o}i's
two-variable normalization, with the curve variable denoted by \(z\),
the NCDT partition function is
\[
  Z_A(q,z)=
  M(-q)^2
  \prod_{k\geq 1}(1-(-q)^kz)^k(1-(-q)^kz^{-1})^k,
\]
while the two commutative resolutions contribute
\[
  Z_{Y_+}(q,z)=M(-q)^2\prod_{k\geq 1}(1-(-q)^kz)^k,
  \qquad
  Z_{Y_-}(q,z)=M(-q)^2\prod_{k\geq 1}(1-(-q)^kz^{-1})^k.
\]
After division by the MacMahon factor, the reduced factorization
\[
  Z_A'(q,z)=Z_{Y_+}'(q,z)\,Z_{Y_-}'(q,z)
\]
is a wall-crossing identity between chambers.  Morrison--Mozgovoy--
Nagao--Szendr\H{o}i make this chamber structure explicit: the negative
quadrant gives NCDT, chambers near \(L_\infty\) give DT and PT for
\(Y_+\), and the opposite side gives the corresponding chambers for
\(Y_-\).

\subsection*{The Conifold CoHA and the \texorpdfstring{\(\kappa_{\mathrm{ch}}\)}{chiral invariant} Value}

The conifold CoHA is the critical CoHA of the Klebanov--Witten
quiver-with-potential:
\[
  \CoHA_{\mathrm{con}}=\CoHA(Q_{\mathrm{con}},W_{\mathrm{KW}}).
\]
Schiffmann--Vasserot's affine-Yangian positive-half theorem is the
one-vertex \(\bC^3\) case:
\[
  \CoHA(\bC^3)=Y^+(\widehat{\mathfrak{gl}}_1).
\]
The conifold calculation uses the two-vertex Jacobi algebra
\(J_{\mathrm{con}}\), the conifold wall-crossing, and the conifold
shuffle kernel.

The resolved conifold
\[
  Y=\operatorname{Tot}\bigl(\cO_{\bP^1}(-1)^{\oplus 2}\bigr)
\]
is a non-compact toric CY$_3$ with one compact curve class.  It lies
outside the local-surface class \(\operatorname{Tot}(K_S)\).  The
local-surface formula
\[
  \kappa_{\mathrm{ch}}=\tfrac{1}{2}\chi_{\mathrm{top}}(S)
\]
belongs to that class of total spaces over Fano surfaces.  The direct
McKay/Jacobi calculation for the conifold gives
\[
  \kappa_{\mathrm{ch}}(A_{\mathrm{con}})=\DT_{(1,0)}=1 .
\]
The local \(\bP^2\) row uses the three-vertex
\(\mathbb Z/3\)-McKay quiver and gives
\(\kappa_{\mathrm{ch}}=3/2\).

The four \(\kappa_\bullet\)-invariants remain separate.  The conifold
statement uses \(\kappa_{\mathrm{ch}}\).  The compact formula
\(\kappa_{\mathrm{cat}}=\chi(\cO_X)\) has no direct non-compact value on
\(Y\); \(\kappa_{\mathrm{BKM}}=c_N(0)/2\) is a Borcherds-product
weight attached to a separate BKM construction; and
\(\kappa_{\mathrm{fiber}}\) belongs to a chosen fibration or lattice
datum.  These invariants require their own constructions.

For this \(d=3\) input, the specialized CY-to-chiral output is
\(E_1\)-chiral.  \(E_2\)-structure belongs to the centre, double, or
module-category layer.
